\documentclass[11pt,twoside]{article}
\usepackage[utf8]{inputenc}
\usepackage[margin=1in]{geometry}
\usepackage{titlesec}
\usepackage{parskip}
\usepackage{fancyhdr}
\usepackage{amsmath, amssymb}
\usepackage[round,sort&compress]{natbib}
\usepackage{tikz}
\usetikzlibrary{positioning, arrows.meta, fit, backgrounds}
\usepackage{hyperref}

% Approximation of JMLR style parameters
\setlength{\parindent}{0pt}
\setlength{\parskip}{0.3em}
\titlespacing*{\section}{0pt}{0.6em}{0.3em}
\titlespacing*{\subsection}{0pt}{0.5em}{0.2em}

\pagestyle{fancy}
\fancyhf{}
\lhead{\footnotesize STAT8101}
\chead{\footnotesize \href{http://www.cs.columbia.edu/~blei/teaching.html}{Invariance, Causality, and Applications}}
\rhead{\footnotesize \href{http://www.cs.columbia.edu/~blei/}{David M. Blei}}
\cfoot{\thepage}

\begin{document}
\vspace*{-1.5cm}
\begin{center}
{\Large \textbf{Reader Report: Week 11}} \\
\vspace{0.2em}
{\large Bayesian Optimization: Regret, Scalability, and the Empirical Bayes Bridge} \\
\vspace{0.3em}
Daniel Hardesty Lewis (\texttt{dl3645}) $\cdot$ Spring 2026
\end{center}
\vspace{-1.5em}
\thispagestyle{fancy}

\section*{Summary}
Building on last week's foundations, I focused on three areas from \citet{garnett2023bayesian}: (1) theoretical regret guarantees for GP-UCB, (2) high-dimensional scaling via random embeddings and trust regions, and (3) the link between marginal likelihood maximization and empirical Bayes.

\textit{Regret bounds.} For the GP-UCB acquisition function $\mu_t(x) + \beta_t^{1/2}\sigma_t(x)$, \citet{srinivas2010ucb} establish a cumulative regret bound $R_T \leq \sqrt{C_1 T \beta_T \gamma_T}$ with high probability, i.e.\ $O^*\!\left(\sqrt{T \gamma_T}\right)$ suppressing log factors in $\beta_T$, where $\gamma_T = \max_{|A|=T} I(\mathbf{y}_A; f)$ is the maximum information gain and $C_1 = 8/\log(1+\sigma^{-2})$. For RBF/SE kernels $\gamma_T = O((\log T)^{d+1})$; for Mat\'{e}rn-$\nu$ kernels $\gamma_T = O(T^{d(d+1)/(2\nu+d(d+1))}(\log T))$ (Theorem 4 + Corollary, arXiv v2 \S5.2), making $\nu$ operationally decisive for convergence speed in spatial domains.

\textit{High-dimensional BO.} Standard GP inference costs $O(n^3)$ in the number of observations $n$, and acquisition function optimization degrades rapidly with input dimensionality $d$; together these render naive BO prohibitive for $d \gg 20$ or large sample budgets. REMBO \citep{wang2013rembo} projects the search space into a low-dimensional random linear subspace, exploiting effective dimensionality in practice. TuRBO \citep{eriksson2019turbo} alternatively maintains local trust regions with independent GP models, escaping the global stationarity assumption that pathologically homogenizes acquisition values across a large domain.

\textit{Marginal likelihood and empirical Bayes.} GP hyperparameters (kernel length-scales, signal variance, noise) are typically estimated by maximizing the log marginal likelihood $\log p(\mathbf{y} \mid X, \theta) = -\tfrac{1}{2}\mathbf{y}^\top (K_\theta + \sigma_n^2 I)^{-1}\mathbf{y} - \tfrac{1}{2}\log|K_\theta + \sigma_n^2 I| - \tfrac{n}{2}\log 2\pi$. This is precisely a Type-II maximum likelihood (empirical Bayes) procedure: treating hyperparameters as fixed point estimates rather than integrating them out. The connection to next week's topic is direct --- empirical Bayes is the mechanism by which BO adapts its prior to the observed function landscape.

\section*{Critique}
The regret analysis reveals a structural tension for causal applications. The bound $O^*(\sqrt{T\gamma_T})$ is derived under a fixed, stationary kernel --- yet in the zoning domain, the mechanism generating housing outcomes shifts after the 2022 municipal election. If the true structural causal model undergoes a discrete change-point, $\gamma_T$ is not a property of a single fixed kernel but of a \textit{sequence} of kernels. No regret bound I am aware of accommodates mid-sequence mechanism shifts of this type; dynamic causal BO \citep{aglietti2021dcbo} addresses gradual distributional drift but not abrupt structural replacement of the causal graph.

The empirical Bayes procedure for kernel hyperparameters is similarly fragile under mechanism shift. Maximizing the marginal likelihood globally averages over pre- and post-shift data, systematically miscalibrating the length-scale. A natural fix --- fitting separate GP models on pre- and post-change-point data --- reintroduces the small-$n$ problem that motivated a pooled GP in the first place. This points toward hierarchical Bayes over change-point location as the appropriate next step, though Garnett does not develop this explicitly.

Finally, the REMBO embedding assumption (that the effective dimensionality of the optimum is low) is plausible for hyperparameter tuning but questionable for spatial policy: geographic spillovers couple dimensions explicitly, so the true intrinsic dimensionality may scale with the number of neighborhoods.

\section*{Questions}
\begin{enumerate}
    \item The GP marginal likelihood is the canonical empirical Bayes objective for BO. But empirical Bayes shrinkage (as in James--Stein or Efron's nonparametric version) operates on a different statistical object --- a prior over many means. Is there a unified framework in which the BO marginal likelihood and Efron-style empirical Bayes are both special cases of the same hierarchical model? Or are they fundamentally different uses of the phrase ``empirical Bayes''?
    \item TuRBO decomposes the global BO problem into local trust regions, each with its own GP. In a causal domain where the action $a$ represents a policy intervention with spillovers (i.e., $f(a_i)$ and $f(a_j)$ are not independent for neighboring regions $i,j$), local GP independence is violated. Are there principled methods for defining trust-region boundaries that respect known causal graph structure rather than Euclidean distance in the action space?
\end{enumerate}

\newpage
\vspace{1em}
\hrule
\vspace{1em}
\section*{Project Update: Simulating the Historical Zoning Sequence as a BO Trajectory}
\thispagestyle{empty}

\begin{abstract}
Last week I proposed framing Austin's 7,074 discretionary zoning cases as a constrained Bayesian Optimization trajectory. This week I began operationalizing that framing: constructing the action space, defining the objective and constraint surrogates, and running a retrospective counterfactual comparing the historical approval sequence against an idealized GP-UCB policy.
\end{abstract}

\section*{Setup and Action Space}
Each zoning case $t$ is an ``evaluation'' of the objective. I define the action $a_t \in \mathcal{A} \subset \mathbb{R}^d$ as a vector of case-level features: proposed density (FAR), parcel area, distance to CBD, zoning category ordinal, and council district indicator. The objective $f(a_t)$ is total permitted housing units (a proxy for yield). The safety constraint $g(a_t) \leq 0$ encodes whether a protest petition was filed; the latent margin $g$ is modeled as a GP as in last week's formulation.

\section*{GP Surrogate Fitting}
I fit separate GP surrogates $f$ and $g$ on the pre-2022 data ($n = 4{,}312$ cases) using a Matérn-$3/2$ kernel with automatic relevance determination (ARD), selecting hyperparameters via marginal likelihood maximization (\emph{empirical Bayes}). Post-2022 cases ($n = 2{,}762$) are held out as a test regime. ARD length-scale estimates from the fitted kernel suggest that three dimensions (density, distance, district) carry most of the explanatory variance --- the remaining dimensions receive near-infinite length-scales, implying near-zero marginal information contribution --- which is consistent with the REMBO hypothesis of a low-dimensional active subspace, though not a rigorous computation of $\gamma_T$.

\section*{Counterfactual Comparison}
I simulate what a GP-UCB policy (with $\beta_t$ scheduled per \citet{srinivas2010ucb}) would have recommended at each decision point, conditioning only on the cases previously approved. Preliminary results suggest the historical council was substantially \textit{more exploitative} than GP-UCB in the pre-2022 period (concentrating approvals in already-dense, low-protest corridors) and then abruptly shifted toward exploration in 2022--2023, consistent with the structural change-point hypothesis. The constraint GP $g$ shows a visible kernel miscalibration across the change-point, supporting last week's critique.

\section*{Next Steps}
Formalize the change-point detection within the GP framework (e.g., by fitting separate GP models on pre- and post-change-point windows, informed by the dynamic causal structure formalism in \citet{aglietti2021dcbo}). Evaluate whether SAFEOPT \citep{sui2015safeopt} --- a single-objective constrained optimizer that guarantees safety with high probability throughout exploration --- would have produced outcomes dominating the historical sequence on both yield and protest rate simultaneously, and quantify the regret gap attributable to the unmodeled structural shift.

\vspace{-0.5em}
{\scriptsize
\setlength{\bibsep}{0pt}
\bibliographystyle{plainnat}
\bibliography{references}
}
\end{document}
